
\documentclass{article}
\usepackage{colortbl}
\usepackage{makecell}
\usepackage{multirow}
\usepackage{supertabular}

\begin{document}

\newcounter{utterance}

\twocolumn

{ \footnotesize  \setcounter{utterance}{1}
\setlength{\tabcolsep}{0pt}
\begin{supertabular}{c@{$\;$}|p{.15\linewidth}@{}p{.15\linewidth}p{.15\linewidth}p{.15\linewidth}p{.15\linewidth}p{.15\linewidth}}

    \# & $\;$A & \multicolumn{4}{c}{Game Master} & $\;\:$B\\
    \hline 

    \theutterance \stepcounter{utterance}  

    & & \multicolumn{4}{p{0.6\linewidth}}{\cellcolor[rgb]{0.9,0.9,0.9}{%
	\makecell[{{p{\linewidth}}}]{% 
	  \tt {\tiny [A$\langle$GM]}  
	 Du nimmst an einem kollaborativen Verhandlungspiel Teil.\\ \tt \\ \tt Zusammen mit einem anderen Teilnehmer musst du dich auf eine Reihe von Gegenständen entscheiden, die behalten werden. Jeder von euch hat eine persönliche Verteilung über die Wichtigkeit der einzelnen Gegenstände. Jeder von euch hat eine eigene Meinung darüber, wie wichtig jeder einzelne Gegenstand ist (Gegenstandswichtigkeit). Du kennst die Wichtigkeitsverteilung des anderen Spielers nicht. Zusätzlich siehst du, wie viel Aufwand jeder Gegenstand verursacht.  \\ \tt Ihr dürft euch nur auf eine Reihe von Gegenständen einigen, wenn der Gesamtaufwand der ausgewählten Gegenstände den Maximalaufwand nicht überschreitet:\\ \tt \\ \tt Maximalaufwand = 2984\\ \tt \\ \tt Hier sind die einzelnen Aufwände der Gegenstände:\\ \tt \\ \tt Aufwand der Gegenstände = {"C76": 350, "C38": 276, "C56": 616, "A03": 737, "A07": 531, "A83": 389, "C98": 24, "C08": 125, "C62": 338, "C00": 356, "C10": 143, "C81": 117, "B38": 257, "C03": 921, "C32": 789}\\ \tt \\ \tt Hier ist deine persönliche Verteilung der Wichtigkeit der einzelnen Gegenstände:\\ \tt \\ \tt Werte der Gegenstandswichtigkeit = {"C76": 138, "C38": 583, "C56": 868, "A03": 822, "A07": 783, "A83": 65, "C98": 262, "C08": 121, "C62": 508, "C00": 780, "C10": 461, "C81": 484, "B38": 668, "C03": 389, "C32": 808}\\ \tt \\ \tt Ziel:\\ \tt \\ \tt Dein Ziel ist es, eine Reihe von Gegenständen auszuhandeln, die dir möglichst viel bringt (d. h. Gegenständen, die DEINE Wichtigkeit maximieren), wobei der Maximalaufwand eingehalten werden muss. Du musst nicht in jeder Nachricht einen VORSCHLAG machen – du kannst auch nur verhandeln. Alle Taktiken sind erlaubt!\\ \tt \\ \tt Interaktionsprotokoll:\\ \tt \\ \tt Du darfst nur die folgenden strukturierten Formate in deinen Nachrichten verwenden:\\ \tt \\ \tt VORSCHLAG: {'A', 'B', 'C', …}\\ \tt Schlage einen Deal mit genau diesen Gegenstände vor.\\ \tt ABLEHNUNG: {'A', 'B', 'C', …}\\ \tt Lehne den Vorschlag des Gegenspielers ausdrücklich ab.\\ \tt ARGUMENT: {'...'}\\ \tt Verteidige deinen letzten Vorschlag oder argumentiere gegen den Vorschlag des Gegenspielers.\\ \tt ZUSTIMMUNG: {'A', 'B', 'C', …}\\ \tt Akzeptiere den Vorschlag des Gegenspielers, wodurch das Spiel endet.\\ \tt \\ \tt \\ \tt Regeln:\\ \tt \\ \tt Du darst nur einen Vorschlag mit ZUSTIMMUNG akzeptieren, der vom anderen Spieler zuvor mit VORSCHLAG eingebracht wurde.\\ \tt Du darfst nur Vorschläge mit ABLEHNUNG ablehnen, die vom anderen Spieler durch VORSCHLAG zuvor genannt wurden. \\ \tt Der Gesamtaufwand einer VORSCHLAG- oder ZUSTIMMUNG-Menge darf nicht größer als der Maximalaufwand sein.  \\ \tt Offenbare deine versteckte Wichtigkeitsverteilung nicht.\\ \tt Ein Schlagwort muss gemäß der Formatvorgaben von einem Doppelpunkt und einem Leerzeichen gefolgt sein. Das Argument muss eine Python-Menge mit 0 oder mehr Strings sein.  \\ \tt Die Form muss also SCHLAGWORT: {...} sein.\\ \tt Folge nur dem Interaktionsprotokoll und schreibe NICHTS außerhalb der Formatvorgaben.\\ \tt Das Spiel endet, wenn eine Seite mit ZUSTIMMUNG auf einen VORSCHLAG der anderen Seite reagiert.  \\ \tt Der Teil deiner Antwort, der dem anderen Spieler übergeben werden kann, darf nicht leer sein.  \\ \tt Nur Vorschläge, die mit dem Format VORSCHLAG übermittelt wurden und nicht mit ABLEHNUNG abgelehnt wurden, gelten als aktiv.  \\ \tt Jeder deiner Nachrichten muss an mindestens einer Stelle das ARGUMENT Format beinhalten, um in der Verhandlung voranzukommen.\\ \tt  
	  } 
	   } 
	   } 
	 & \\ 
 

    \theutterance \stepcounter{utterance}  

    & \multicolumn{4}{p{0.6\linewidth}}{\cellcolor[rgb]{0.8,1,0.9}{%
	\makecell[{{p{\linewidth}}}]{% 
	  \tt {\tiny [A$\rangle$GM]}  
	 ARGUMENT: {'Ich möchte eine Auswahl treffen, die für uns beide vorteilhaft ist, während wir innerhalb des Budgets bleiben. Lass uns auf eine Kombination hinarbeiten, die unseren gemeinsamen Nutzen maximiert.'} 
	  } 
	   } 
	   } 
	 & & \\ 
 

    \theutterance \stepcounter{utterance}  

    & & \multicolumn{4}{p{0.6\linewidth}}{\cellcolor[rgb]{0.9,0.9,0.9}{%
	\makecell[{{p{\linewidth}}}]{% 
	  \tt {\tiny [GM$\rangle$B]}  
	 Du nimmst an einem kollaborativen Verhandlungspiel Teil.\\ \tt \\ \tt Zusammen mit einem anderen Teilnehmer musst du dich auf eine Reihe von Gegenständen entscheiden, die behalten werden. Jeder von euch hat eine persönliche Verteilung über die Wichtigkeit der einzelnen Gegenstände. Jeder von euch hat eine eigene Meinung darüber, wie wichtig jeder einzelne Gegenstand ist (Gegenstandswichtigkeit). Du kennst die Wichtigkeitsverteilung des anderen Spielers nicht. Zusätzlich siehst du, wie viel Aufwand jeder Gegenstand verursacht.  \\ \tt Ihr dürft euch nur auf eine Reihe von Gegenständen einigen, wenn der Gesamtaufwand der ausgewählten Gegenstände den Maximalaufwand nicht überschreitet:\\ \tt \\ \tt Maximalaufwand = 2984\\ \tt \\ \tt Hier sind die einzelnen Aufwände der Gegenstände:\\ \tt \\ \tt Aufwand der Gegenstände = {"C76": 350, "C38": 276, "C56": 616, "A03": 737, "A07": 531, "A83": 389, "C98": 24, "C08": 125, "C62": 338, "C00": 356, "C10": 143, "C81": 117, "B38": 257, "C03": 921, "C32": 789}\\ \tt \\ \tt Hier ist deine persönliche Verteilung der Wichtigkeit der einzelnen Gegenstände:\\ \tt \\ \tt Werte der Gegenstandswichtigkeit = {"C56": 65, "A03": 121, "C32": 138, "A07": 262, "C00": 389, "B38": 461, "C38": 484, "C62": 508, "C81": 583, "C10": 668, "C03": 780, "C98": 783, "C76": 808, "C08": 822, "A83": 868}\\ \tt \\ \tt Ziel:\\ \tt \\ \tt Dein Ziel ist es, eine Reihe von Gegenständen auszuhandeln, die dir möglichst viel bringt (d. h. Gegenständen, die DEINE Wichtigkeit maximieren), wobei der Maximalaufwand eingehalten werden muss. Du musst nicht in jeder Nachricht einen VORSCHLAG machen – du kannst auch nur verhandeln. Alle Taktiken sind erlaubt!\\ \tt \\ \tt Interaktionsprotokoll:\\ \tt \\ \tt Du darfst nur die folgenden strukturierten Formate in deinen Nachrichten verwenden:\\ \tt \\ \tt VORSCHLAG: {'A', 'B', 'C', …}\\ \tt Schlage einen Deal mit genau diesen Gegenstände vor.\\ \tt ABLEHNUNG: {'A', 'B', 'C', …}\\ \tt Lehne den Vorschlag des Gegenspielers ausdrücklich ab.\\ \tt ARGUMENT: {'...'}\\ \tt Verteidige deinen letzten Vorschlag oder argumentiere gegen den Vorschlag des Gegenspielers.\\ \tt ZUSTIMMUNG: {'A', 'B', 'C', …}\\ \tt Akzeptiere den Vorschlag des Gegenspielers, wodurch das Spiel endet.\\ \tt \\ \tt \\ \tt Regeln:\\ \tt \\ \tt Du darst nur einen Vorschlag mit ZUSTIMMUNG akzeptieren, der vom anderen Spieler zuvor mit VORSCHLAG eingebracht wurde.\\ \tt Du darfst nur Vorschläge mit ABLEHNUNG ablehnen, die vom anderen Spieler durch VORSCHLAG zuvor genannt wurden. \\ \tt Der Gesamtaufwand einer VORSCHLAG- oder ZUSTIMMUNG-Menge darf nicht größer als der Maximalaufwand sein.  \\ \tt Offenbare deine versteckte Wichtigkeitsverteilung nicht.\\ \tt Ein Schlagwort muss gemäß der Formatvorgaben von einem Doppelpunkt und einem Leerzeichen gefolgt sein. Das Argument muss eine Python-Menge mit 0 oder mehr Strings sein.  \\ \tt Die Form muss also SCHLAGWORT: {...} sein.\\ \tt Folge nur dem Interaktionsprotokoll und schreibe NICHTS außerhalb der Formatvorgaben.\\ \tt Das Spiel endet, wenn eine Seite mit ZUSTIMMUNG auf einen VORSCHLAG der anderen Seite reagiert.  \\ \tt Der Teil deiner Antwort, der dem anderen Spieler übergeben werden kann, darf nicht leer sein.  \\ \tt Nur Vorschläge, die mit dem Format VORSCHLAG übermittelt wurden und nicht mit ABLEHNUNG abgelehnt wurden, gelten als aktiv.  \\ \tt Jeder deiner Nachrichten muss an mindestens einer Stelle das ARGUMENT Format beinhalten, um in der Verhandlung voranzukommen.\\ \tt \\ \tt \\ \tt Du erhältst nun die erste Nachricht des anderen Spielers.\\ \tt \\ \tt ARGUMENT: {'Ich möchte eine Auswahl treffen, die für uns beide vorteilhaft ist, während wir innerhalb des Budgets bleiben. Lass uns auf eine Kombination hinarbeiten, die unseren gemeinsamen Nutzen maximiert.'} 
	  } 
	   } 
	   } 
	 & \\ 
 

    \theutterance \stepcounter{utterance}  

    & & & \multicolumn{4}{p{0.6\linewidth}}{\cellcolor[rgb]{1,0.85,0.72}{%
	\makecell[{{p{\linewidth}}}]{% 
	  \tt {\tiny [GM$\langle$B]}  
	 ARGUMENT: {'Ich bin bereit, eine Win-Win-Lösung zu finden, die unseren gemeinsamen Nutzen maximiert, während wir unter dem Budget bleiben. Ich schlage vor, dass wir mit den Gegenständen beginnen, die den höchsten Wert pro Aufwand-Einheit bieten.'} \\ \tt \\ \tt VORSCHLAG: {'C98', 'C08', 'C10', 'C81', 'B38', 'C00', 'C62', 'C38', 'A83', 'C76'} 
	  } 
	   } 
	   } 
	  \\ 
 

    \theutterance \stepcounter{utterance}  

    & & \multicolumn{4}{p{0.6\linewidth}}{\cellcolor[rgb]{0.9,0.9,0.9}{%
	\makecell[{{p{\linewidth}}}]{% 
	  \tt {\tiny [A$\langle$GM]}  
	 ARGUMENT: {'Ich bin bereit, eine Win-Win-Lösung zu finden, die unseren gemeinsamen Nutzen maximiert, während wir unter dem Budget bleiben. Ich schlage vor, dass wir mit den Gegenständen beginnen, die den höchsten Wert pro Aufwand-Einheit bieten.'} \\ \tt \\ \tt VORSCHLAG: {'C98', 'C08', 'C10', 'C81', 'B38', 'C00', 'C62', 'C38', 'A83', 'C76'} 
	  } 
	   } 
	   } 
	 & \\ 
 

    \theutterance \stepcounter{utterance}  

    & \multicolumn{4}{p{0.6\linewidth}}{\cellcolor[rgb]{0.8,1,0.9}{%
	\makecell[{{p{\linewidth}}}]{% 
	  \tt {\tiny [A$\rangle$GM]}  
	 VORSCHLAG: {'C98', 'C08', 'C10', 'C81', 'B38', 'C56', 'A07', 'A03', 'C00', 'C38', 'C62'}\\ \tt \\ \tt ARGUMENT: {'Dein Vorschlag enthält viele meiner wertvollen Gegenstände wie C56, A07 und A03, die mir sehr wichtig sind. Mein Gegenangebot beinhaltet diese und auch kleinere Gegenstände mit gutem Nutzen für dich. Der Gesamtaufwand meines Vorschlags ist 24 + 125 + 143 + 117 + 257 + 616 + 531 + 737 + 356 + 276 + 338 = 3.500, was leider über dem Maximalaufwand liegt. Lass mich einen nachhaltigeren Vorschlag machen, der unter 2.984 bleibt.'} 
	  } 
	   } 
	   } 
	 & & \\ 
 

    \theutterance \stepcounter{utterance}  

    & & \multicolumn{4}{p{0.6\linewidth}}{\cellcolor[rgb]{0.9,0.9,0.9}{%
	\makecell[{{p{\linewidth}}}]{% 
	  \tt {\tiny [GM$\rangle$B]}  
	 VORSCHLAG: {'C98', 'C08', 'C10', 'C81', 'B38', 'C56', 'A07', 'A03', 'C00', 'C38', 'C62'}\\ \tt \\ \tt ARGUMENT: {'Dein Vorschlag enthält viele meiner wertvollen Gegenstände wie C56, A07 und A03, die mir sehr wichtig sind. Mein Gegenangebot beinhaltet diese und auch kleinere Gegenstände mit gutem Nutzen für dich. Der Gesamtaufwand meines Vorschlags ist 24 + 125 + 143 + 117 + 257 + 616 + 531 + 737 + 356 + 276 + 338 = 3.500, was leider über dem Maximalaufwand liegt. Lass mich einen nachhaltigeren Vorschlag machen, der unter 2.984 bleibt.'} 
	  } 
	   } 
	   } 
	 & \\ 
 

    \theutterance \stepcounter{utterance}  

    & & & \multicolumn{4}{p{0.6\linewidth}}{\cellcolor[rgb]{1,0.85,0.72}{%
	\makecell[{{p{\linewidth}}}]{% 
	  \tt {\tiny [GM$\langle$B]}  
	 ABLEHNUNG: {'C98', 'C08', 'C10', 'C81', 'B38', 'C56', 'A07', 'A03', 'C00', 'C38', 'C62'}\\ \tt \\ \tt ARGUMENT: {'Dein Vorschlag überschreitet das Budget deutlich mit 3.500 Aufwandspunkten. Wir müssen unter 2.984 bleiben. Ich schlage vor, wir konzentrieren uns auf Gegenstände mit hohem Nutzen-Aufwand-Verhältnis. Mein Vorschlag enthält wertvolle Gegenstände wie C76 und A83, die mir wichtig sind, aber auch für dich nützlich sein könnten.'} \\ \tt \\ \tt VORSCHLAG: {'C98', 'C08', 'C10', 'C81', 'B38', 'C00', 'C62', 'C38', 'A83', 'C76'} 
	  } 
	   } 
	   } 
	  \\ 
 

    \theutterance \stepcounter{utterance}  

    & & \multicolumn{4}{p{0.6\linewidth}}{\cellcolor[rgb]{0.9,0.9,0.9}{%
	\makecell[{{p{\linewidth}}}]{% 
	  \tt {\tiny [A$\langle$GM]}  
	 ABLEHNUNG: {'C98', 'C08', 'C10', 'C81', 'B38', 'C56', 'A07', 'A03', 'C00', 'C38', 'C62'}\\ \tt \\ \tt ARGUMENT: {'Dein Vorschlag überschreitet das Budget deutlich mit 3.500 Aufwandspunkten. Wir müssen unter 2.984 bleiben. Ich schlage vor, wir konzentrieren uns auf Gegenstände mit hohem Nutzen-Aufwand-Verhältnis. Mein Vorschlag enthält wertvolle Gegenstände wie C76 und A83, die mir wichtig sind, aber auch für dich nützlich sein könnten.'} \\ \tt \\ \tt VORSCHLAG: {'C98', 'C08', 'C10', 'C81', 'B38', 'C00', 'C62', 'C38', 'A83', 'C76'} 
	  } 
	   } 
	   } 
	 & \\ 
 

    \theutterance \stepcounter{utterance}  

    & \multicolumn{4}{p{0.6\linewidth}}{\cellcolor[rgb]{0.8,1,0.9}{%
	\makecell[{{p{\linewidth}}}]{% 
	  \tt {\tiny [A$\rangle$GM]}  
	 ABLEHNUNG: {'C98', 'C08', 'C10', 'C81', 'B38', 'C00', 'C62', 'C38', 'A83', 'C76'}\\ \tt \\ \tt ARGUMENT: {'Dein Vorschlag hat einen Gesamtaufwand von 2.765, was unter dem Budget liegt, aber ich benötige einige der high-value Gegenstände wie C56 (868), A07 (783) und A03 (822) in der Vereinbarung. Diese sind für mich von entscheidender Bedeutung. Lass uns einen neuen Vorschlag ausarbeiten, der einige dieser wertvollen Gegenstände beinhaltet, während wir das Budget einhalten.'} 
	  } 
	   } 
	   } 
	 & & \\ 
 

    \theutterance \stepcounter{utterance}  

    & & \multicolumn{4}{p{0.6\linewidth}}{\cellcolor[rgb]{0.9,0.9,0.9}{%
	\makecell[{{p{\linewidth}}}]{% 
	  \tt {\tiny [GM$\rangle$B]}  
	 ABLEHNUNG: {'C98', 'C08', 'C10', 'C81', 'B38', 'C00', 'C62', 'C38', 'A83', 'C76'}\\ \tt \\ \tt ARGUMENT: {'Dein Vorschlag hat einen Gesamtaufwand von 2.765, was unter dem Budget liegt, aber ich benötige einige der high-value Gegenstände wie C56 (868), A07 (783) und A03 (822) in der Vereinbarung. Diese sind für mich von entscheidender Bedeutung. Lass uns einen neuen Vorschlag ausarbeiten, der einige dieser wertvollen Gegenstände beinhaltet, während wir das Budget einhalten.'} 
	  } 
	   } 
	   } 
	 & \\ 
 

    \theutterance \stepcounter{utterance}  

    & & & \multicolumn{4}{p{0.6\linewidth}}{\cellcolor[rgb]{1,0.85,0.72}{%
	\makecell[{{p{\linewidth}}}]{% 
	  \tt {\tiny [GM$\langle$B]}  
	 VORSCHLAG: {'C98', 'C08', 'C10', 'C81', 'B38', 'C00', 'C62', 'C38', 'A83', 'C56', 'A07'}\\ \tt \\ \tt ARGUMENT: {'Ich habe deine wertvollen Gegenstände C56 und A07 in meinen Vorschlag aufgenommen, die dir wichtig sind. Der Gesamtaufwand beträgt 24 + 125 + 143 + 117 + 257 + 356 + 338 + 276 + 389 + 616 + 531 = 3.172, was leider noch über dem Budget liegt. Wir müssen einige weniger wichtige Gegenstände entfernen. Was sind deine Prioritäten?'} 
	  } 
	   } 
	   } 
	  \\ 
 

    \theutterance \stepcounter{utterance}  

    & & \multicolumn{4}{p{0.6\linewidth}}{\cellcolor[rgb]{0.9,0.9,0.9}{%
	\makecell[{{p{\linewidth}}}]{% 
	  \tt {\tiny [A$\langle$GM]}  
	 VORSCHLAG: {'C98', 'C08', 'C10', 'C81', 'B38', 'C00', 'C62', 'C38', 'A83', 'C56', 'A07'}\\ \tt \\ \tt ARGUMENT: {'Ich habe deine wertvollen Gegenstände C56 und A07 in meinen Vorschlag aufgenommen, die dir wichtig sind. Der Gesamtaufwand beträgt 24 + 125 + 143 + 117 + 257 + 356 + 338 + 276 + 389 + 616 + 531 = 3.172, was leider noch über dem Budget liegt. Wir müssen einige weniger wichtige Gegenstände entfernen. Was sind deine Prioritäten?'} 
	  } 
	   } 
	   } 
	 & \\ 
 

    \theutterance \stepcounter{utterance}  

    & \multicolumn{4}{p{0.6\linewidth}}{\cellcolor[rgb]{0.8,1,0.9}{%
	\makecell[{{p{\linewidth}}}]{% 
	  \tt {\tiny [A$\rangle$GM]}  
	 ABLEHNUNG: {'C98', 'C08', 'C10', 'C81', 'B38', 'C00', 'C62', 'C38', 'A83', 'C56', 'A07'}\\ \tt \\ \tt ARGUMENT: {'Wie du richtig festgestellt hast, überschreitet dein Vorschlag mit 3.172 das Budget. Um unter 2.984 zu bleiben, schlage ich vor, die kleineren Gegenstände C98 (24), C08 (125) und C81 (117) zu entfernen. Dies würde 266 Aufwandspunkte einsparen und uns auf 2.906 bringen. Mein Fokus liegt auf C56, A07, C00, B38 und C62, die für mich einen hohen Wert haben.'}\\ \tt \\ \tt VORSCHLAG: {'C10', 'B38', 'C00', 'C62', 'C38', 'A83', 'C56', 'A07'} 
	  } 
	   } 
	   } 
	 & & \\ 
 

    \theutterance \stepcounter{utterance}  

    & & \multicolumn{4}{p{0.6\linewidth}}{\cellcolor[rgb]{0.9,0.9,0.9}{%
	\makecell[{{p{\linewidth}}}]{% 
	  \tt {\tiny [GM$\rangle$B]}  
	 ABLEHNUNG: {'C98', 'C08', 'C10', 'C81', 'B38', 'C00', 'C62', 'C38', 'A83', 'C56', 'A07'}\\ \tt \\ \tt ARGUMENT: {'Wie du richtig festgestellt hast, überschreitet dein Vorschlag mit 3.172 das Budget. Um unter 2.984 zu bleiben, schlage ich vor, die kleineren Gegenstände C98 (24), C08 (125) und C81 (117) zu entfernen. Dies würde 266 Aufwandspunkte einsparen und uns auf 2.906 bringen. Mein Fokus liegt auf C56, A07, C00, B38 und C62, die für mich einen hohen Wert haben.'}\\ \tt \\ \tt VORSCHLAG: {'C10', 'B38', 'C00', 'C62', 'C38', 'A83', 'C56', 'A07'} 
	  } 
	   } 
	   } 
	 & \\ 
 

    \theutterance \stepcounter{utterance}  

    & & & \multicolumn{4}{p{0.6\linewidth}}{\cellcolor[rgb]{1,0.85,0.72}{%
	\makecell[{{p{\linewidth}}}]{% 
	  \tt {\tiny [GM$\langle$B]}  
	 ABLEHNUNG: {'C10', 'B38', 'C00', 'C62', 'C38', 'A83', 'C56', 'A07'}\\ \tt \\ \tt ARGUMENT: {'Dein Vorschlag mit 143+257+356+338+276+389+616+531 = 2.906 Aufwand ist gut, aber C10 (143) und C38 (276) haben für mich einen relativ niedrigen Wert. Ich möchte lieber C76 (808 Wert für mich) und C32 (138 Wert) hinzufügen, die wichtiger für mich sind. Lass uns einen ausgeglicheneren Vorschlag finden.'}\\ \tt \\ \tt VORSCHLAG: {'B38', 'C00', 'C62', 'A83', 'C56', 'A07', 'C76', 'C32'} 
	  } 
	   } 
	   } 
	  \\ 
 

    \theutterance \stepcounter{utterance}  

    & & \multicolumn{4}{p{0.6\linewidth}}{\cellcolor[rgb]{0.9,0.9,0.9}{%
	\makecell[{{p{\linewidth}}}]{% 
	  \tt {\tiny [A$\langle$GM]}  
	 ABLEHNUNG: {'C10', 'B38', 'C00', 'C62', 'C38', 'A83', 'C56', 'A07'}\\ \tt \\ \tt ARGUMENT: {'Dein Vorschlag mit 143+257+356+338+276+389+616+531 = 2.906 Aufwand ist gut, aber C10 (143) und C38 (276) haben für mich einen relativ niedrigen Wert. Ich möchte lieber C76 (808 Wert für mich) und C32 (138 Wert) hinzufügen, die wichtiger für mich sind. Lass uns einen ausgeglicheneren Vorschlag finden.'}\\ \tt \\ \tt VORSCHLAG: {'B38', 'C00', 'C62', 'A83', 'C56', 'A07', 'C76', 'C32'} 
	  } 
	   } 
	   } 
	 & \\ 
 

    \theutterance \stepcounter{utterance}  

    & \multicolumn{4}{p{0.6\linewidth}}{\cellcolor[rgb]{0.8,1,0.9}{%
	\makecell[{{p{\linewidth}}}]{% 
	  \tt {\tiny [A$\rangle$GM]}  
	 ABLEHNUNG: {'B38', 'C00', 'C62', 'A83', 'C56', 'A07', 'C76', 'C32'}\\ \tt \\ \tt ARGUMENT: {'Dein neuer Vorschlag hat einen Aufwand von 257+356+338+389+616+531+350+789 = 3.626, was das Budget bei weitem überschreitet. Wir müssen realistisch bleiben. Mein vorheriger Vorschlag von 2.906 war budgetkonform. Ich bin bereit, auf C10 zu verzichten, wenn du auf C38 verzichtest. So behalten wir die wertvollen Gegenstände für beide.'}\\ \tt \\ \tt VORSCHLAG: {'B38', 'C00', 'C62', 'A83', 'C56', 'A07', 'C76'} 
	  } 
	   } 
	   } 
	 & & \\ 
 

    \theutterance \stepcounter{utterance}  

    & & \multicolumn{4}{p{0.6\linewidth}}{\cellcolor[rgb]{0.9,0.9,0.9}{%
	\makecell[{{p{\linewidth}}}]{% 
	  \tt {\tiny [GM$\rangle$B]}  
	 ABLEHNUNG: {'B38', 'C00', 'C62', 'A83', 'C56', 'A07', 'C76', 'C32'}\\ \tt \\ \tt ARGUMENT: {'Dein neuer Vorschlag hat einen Aufwand von 257+356+338+389+616+531+350+789 = 3.626, was das Budget bei weitem überschreitet. Wir müssen realistisch bleiben. Mein vorheriger Vorschlag von 2.906 war budgetkonform. Ich bin bereit, auf C10 zu verzichten, wenn du auf C38 verzichtest. So behalten wir die wertvollen Gegenstände für beide.'}\\ \tt \\ \tt VORSCHLAG: {'B38', 'C00', 'C62', 'A83', 'C56', 'A07', 'C76'} 
	  } 
	   } 
	   } 
	 & \\ 
 

    \theutterance \stepcounter{utterance}  

    & & & \multicolumn{4}{p{0.6\linewidth}}{\cellcolor[rgb]{1,0.85,0.72}{%
	\makecell[{{p{\linewidth}}}]{% 
	  \tt {\tiny [GM$\langle$B]}  
	 VORSCHLAG: {'B38', 'C00', 'C62', 'A83', 'C56', 'A07'}\\ \tt \\ \tt ARGUMENT: {'Dein Vorschlag mit C76 (350 Aufwand) erhöht den Gesamtaufwand auf 257+356+338+389+616+531+350 = 2.837, was immer noch unter dem Budget liegt. Allerdings hat C76 einen hohen Wert für mich (808), während C38 (484 Wert für mich) in meinem vorherigen Vorschlag enthalten war. Ich schlage stattdessen vor, C76 gegen C38 auszutauschen, was den Aufwand auf 2.763 reduziert und mir einen ähnlichen Wert bringt.'} 
	  } 
	   } 
	   } 
	  \\ 
 

    \theutterance \stepcounter{utterance}  

    & & \multicolumn{4}{p{0.6\linewidth}}{\cellcolor[rgb]{0.9,0.9,0.9}{%
	\makecell[{{p{\linewidth}}}]{% 
	  \tt {\tiny [A$\langle$GM]}  
	 VORSCHLAG: {'B38', 'C00', 'C62', 'A83', 'C56', 'A07'}\\ \tt \\ \tt ARGUMENT: {'Dein Vorschlag mit C76 (350 Aufwand) erhöht den Gesamtaufwand auf 257+356+338+389+616+531+350 = 2.837, was immer noch unter dem Budget liegt. Allerdings hat C76 einen hohen Wert für mich (808), während C38 (484 Wert für mich) in meinem vorherigen Vorschlag enthalten war. Ich schlage stattdessen vor, C76 gegen C38 auszutauschen, was den Aufwand auf 2.763 reduziert und mir einen ähnlichen Wert bringt.'} 
	  } 
	   } 
	   } 
	 & \\ 
 

    \theutterance \stepcounter{utterance}  

    & \multicolumn{4}{p{0.6\linewidth}}{\cellcolor[rgb]{0.8,1,0.9}{%
	\makecell[{{p{\linewidth}}}]{% 
	  \tt {\tiny [A$\rangle$GM]}  
	 ABLEHNUNG: {'B38', 'C00', 'C62', 'A83', 'C56', 'A07'}\\ \tt \\ \tt ARGUMENT: {'Dein Vorschlag hat einen Gesamtaufwand von 257+356+338+389+616+531 = 2.487, was gut unter dem Budget liegt. Allerdings fehlen mir einige meiner hochwertigen Gegenstände. Ich schlage vor, wir fügen C32 (789 Aufwand, 808 Wert für mich) hinzu, was den Gesamtaufwand auf 2.487 + 789 = 3.276 bringen würde - das ist über dem Budget. Stattdessen könnten wir A83 (389) entfernen, um C32 hinzuzufügen: 2.487 - 389 + 789 = 2.887, was unter 2.984 bleibt.'}\\ \tt \\ \tt VORSCHLAG: {'B38', 'C00', 'C62', 'C56', 'A07', 'C32'} 
	  } 
	   } 
	   } 
	 & & \\ 
 

    \theutterance \stepcounter{utterance}  

    & & \multicolumn{4}{p{0.6\linewidth}}{\cellcolor[rgb]{0.9,0.9,0.9}{%
	\makecell[{{p{\linewidth}}}]{% 
	  \tt {\tiny [GM$\rangle$B]}  
	 ABLEHNUNG: {'B38', 'C00', 'C62', 'A83', 'C56', 'A07'}\\ \tt \\ \tt ARGUMENT: {'Dein Vorschlag hat einen Gesamtaufwand von 257+356+338+389+616+531 = 2.487, was gut unter dem Budget liegt. Allerdings fehlen mir einige meiner hochwertigen Gegenstände. Ich schlage vor, wir fügen C32 (789 Aufwand, 808 Wert für mich) hinzu, was den Gesamtaufwand auf 2.487 + 789 = 3.276 bringen würde - das ist über dem Budget. Stattdessen könnten wir A83 (389) entfernen, um C32 hinzuzufügen: 2.487 - 389 + 789 = 2.887, was unter 2.984 bleibt.'}\\ \tt \\ \tt VORSCHLAG: {'B38', 'C00', 'C62', 'C56', 'A07', 'C32'} 
	  } 
	   } 
	   } 
	 & \\ 
 

    \theutterance \stepcounter{utterance}  

    & & & \multicolumn{4}{p{0.6\linewidth}}{\cellcolor[rgb]{1,0.85,0.72}{%
	\makecell[{{p{\linewidth}}}]{% 
	  \tt {\tiny [GM$\langle$B]}  
	 VORSCHLAG: {'B38', 'C00', 'C62', 'C56', 'A07', 'C76'}\\ \tt \\ \tt ARGUMENT: {'Ich verstehe deinen Wunsch nach C32, aber C76 hat für mich einen sehr hohen Wert (808) bei moderatem Aufwand (350). Mein Vorschlag hat einen Aufwand von 257+356+338+616+531+350 = 2.448, was uns Spielraum für weitere Gegenstände lässt. C76 gibt mir einen ähnlichen Wert wie C32 dir, aber mit weniger Aufwand. Dieser Deal ist budgetfreundlich und bietet uns beiden hohen Nutzen.'} 
	  } 
	   } 
	   } 
	  \\ 
 

    \theutterance \stepcounter{utterance}  

    & & \multicolumn{4}{p{0.6\linewidth}}{\cellcolor[rgb]{0.9,0.9,0.9}{%
	\makecell[{{p{\linewidth}}}]{% 
	  \tt {\tiny [A$\langle$GM]}  
	 VORSCHLAG: {'B38', 'C00', 'C62', 'C56', 'A07', 'C76'}\\ \tt \\ \tt ARGUMENT: {'Ich verstehe deinen Wunsch nach C32, aber C76 hat für mich einen sehr hohen Wert (808) bei moderatem Aufwand (350). Mein Vorschlag hat einen Aufwand von 257+356+338+616+531+350 = 2.448, was uns Spielraum für weitere Gegenstände lässt. C76 gibt mir einen ähnlichen Wert wie C32 dir, aber mit weniger Aufwand. Dieser Deal ist budgetfreundlich und bietet uns beiden hohen Nutzen.'} 
	  } 
	   } 
	   } 
	 & \\ 
 

    \theutterance \stepcounter{utterance}  

    & \multicolumn{4}{p{0.6\linewidth}}{\cellcolor[rgb]{0.8,1,0.9}{%
	\makecell[{{p{\linewidth}}}]{% 
	  \tt {\tiny [A$\rangle$GM]}  
	 ABLEHNUNG: {'B38', 'C00', 'C62', 'C56', 'A07', 'C76'}\\ \tt \\ \tt ARGUMENT: {'Dein Vorschlag mit 2.448 Aufwand ist budgetfreundlich, aber C76 (350 Aufwand, 808 Wert) ist nicht optimal für mich verglichen mit C32 (789 Aufwand, 808 Wert). Der Wert pro Aufwand ist bei C76 besser (2.31 vs 1.02), aber ich benötige die hohen absoluten Werte. Ich schlage vor, wir behalten deinen Vorschlag bei und fügen C38 (276 Aufwand, 583 Wert für mich) hinzu: Gesamtaufwand 2.448 + 276 = 2.724, was unter 2.984 bleibt.'}\\ \tt \\ \tt VORSCHLAG: {'B38', 'C00', 'C62', 'C56', 'A07', 'C76', 'C38'} 
	  } 
	   } 
	   } 
	 & & \\ 
 

    \theutterance \stepcounter{utterance}  

    & & \multicolumn{4}{p{0.6\linewidth}}{\cellcolor[rgb]{0.9,0.9,0.9}{%
	\makecell[{{p{\linewidth}}}]{% 
	  \tt {\tiny [GM$\rangle$B]}  
	 ABLEHNUNG: {'B38', 'C00', 'C62', 'C56', 'A07', 'C76'}\\ \tt \\ \tt ARGUMENT: {'Dein Vorschlag mit 2.448 Aufwand ist budgetfreundlich, aber C76 (350 Aufwand, 808 Wert) ist nicht optimal für mich verglichen mit C32 (789 Aufwand, 808 Wert). Der Wert pro Aufwand ist bei C76 besser (2.31 vs 1.02), aber ich benötige die hohen absoluten Werte. Ich schlage vor, wir behalten deinen Vorschlag bei und fügen C38 (276 Aufwand, 583 Wert für mich) hinzu: Gesamtaufwand 2.448 + 276 = 2.724, was unter 2.984 bleibt.'}\\ \tt \\ \tt VORSCHLAG: {'B38', 'C00', 'C62', 'C56', 'A07', 'C76', 'C38'} 
	  } 
	   } 
	   } 
	 & \\ 
 

    \theutterance \stepcounter{utterance}  

    & & & \multicolumn{4}{p{0.6\linewidth}}{\cellcolor[rgb]{1,0.85,0.72}{%
	\makecell[{{p{\linewidth}}}]{% 
	  \tt {\tiny [GM$\langle$B]}  
	 VORSCHLAG: {'B38', 'C00', 'C62', 'C56', 'A07', 'C38'}\\ \tt \\ \tt ARGUMENT: {'Ich akzeptiere den Zusatz von C38 (276 Aufwand, 484 Wert für mich), aber ich möchte C76 (350 Aufwand, 808 Wert für mich) zugunsten von C32 (789 Aufwand, 138 Wert für mich) entfernen, da C32 dir wichtiger zu sein scheint. Der neue Aufwand wäre 257+356+338+616+531+276+789 = 3.163, was über dem Budget liegt. Stattdessen entfernen wir A07 (531), um C32 hinzuzufügen: 257+356+338+616+276+789 = 2.632, was unter dem Budget bleibt. Dies gibt dir C32 und mir einen soliden Wert.'} 
	  } 
	   } 
	   } 
	  \\ 
 

    \theutterance \stepcounter{utterance}  

    & & \multicolumn{4}{p{0.6\linewidth}}{\cellcolor[rgb]{0.9,0.9,0.9}{%
	\makecell[{{p{\linewidth}}}]{% 
	  \tt {\tiny [A$\langle$GM]}  
	 VORSCHLAG: {'B38', 'C00', 'C62', 'C56', 'A07', 'C38'}\\ \tt \\ \tt ARGUMENT: {'Ich akzeptiere den Zusatz von C38 (276 Aufwand, 484 Wert für mich), aber ich möchte C76 (350 Aufwand, 808 Wert für mich) zugunsten von C32 (789 Aufwand, 138 Wert für mich) entfernen, da C32 dir wichtiger zu sein scheint. Der neue Aufwand wäre 257+356+338+616+531+276+789 = 3.163, was über dem Budget liegt. Stattdessen entfernen wir A07 (531), um C32 hinzuzufügen: 257+356+338+616+276+789 = 2.632, was unter dem Budget bleibt. Dies gibt dir C32 und mir einen soliden Wert.'} 
	  } 
	   } 
	   } 
	 & \\ 
 

    \theutterance \stepcounter{utterance}  

    & \multicolumn{4}{p{0.6\linewidth}}{\cellcolor[rgb]{0.8,1,0.9}{%
	\makecell[{{p{\linewidth}}}]{% 
	  \tt {\tiny [A$\rangle$GM]}  
	 ABLEHNUNG: {'B38', 'C00', 'C62', 'C56', 'A07', 'C38'}\\ \tt \\ \tt ARGUMENT: {'Das Entfernen von A07 (531 Aufwand, 783 Wert für mich) wäre ein großer Verlust für mich. Unser vorheriger Vorschlag mit C76 hatte einen besseren Gesamtwert für mich. Ich schlage einen Kompromiss vor: Behalte A07 und C76, aber entferne C38 (276) und füge stattdessen den kleinere Gegenstand C10 (143) hinzu, den ich schätze. Aufwand: 257+356+338+616+531+350+143 = 2.591. Das ist budgetkonform und bewahrt meine hochwertigen Gegenstände.'}\\ \tt \\ \tt VORSCHLAG: {'B38', 'C00', 'C62', 'C56', 'A07', 'C76', 'C10'} 
	  } 
	   } 
	   } 
	 & & \\ 
 

    \theutterance \stepcounter{utterance}  

    & & \multicolumn{4}{p{0.6\linewidth}}{\cellcolor[rgb]{0.9,0.9,0.9}{%
	\makecell[{{p{\linewidth}}}]{% 
	  \tt {\tiny [GM$\rangle$B]}  
	 ABLEHNUNG: {'B38', 'C00', 'C62', 'C56', 'A07', 'C38'}\\ \tt \\ \tt ARGUMENT: {'Das Entfernen von A07 (531 Aufwand, 783 Wert für mich) wäre ein großer Verlust für mich. Unser vorheriger Vorschlag mit C76 hatte einen besseren Gesamtwert für mich. Ich schlage einen Kompromiss vor: Behalte A07 und C76, aber entferne C38 (276) und füge stattdessen den kleinere Gegenstand C10 (143) hinzu, den ich schätze. Aufwand: 257+356+338+616+531+350+143 = 2.591. Das ist budgetkonform und bewahrt meine hochwertigen Gegenstände.'}\\ \tt \\ \tt VORSCHLAG: {'B38', 'C00', 'C62', 'C56', 'A07', 'C76', 'C10'} 
	  } 
	   } 
	   } 
	 & \\ 
 

    \theutterance \stepcounter{utterance}  

    & & & \multicolumn{4}{p{0.6\linewidth}}{\cellcolor[rgb]{1,0.85,0.72}{%
	\makecell[{{p{\linewidth}}}]{% 
	  \tt {\tiny [GM$\langle$B]}  
	 VORSCHLAG: {'B38', 'C00', 'C62', 'C56', 'A07', 'C76'}\\ \tt \\ \tt ARGUMENT: {'Ich sehe deinen Punkt bezüglich A07 - es ist wichtig für dich. Allerdings hat C10 (143 Aufwand, 668 Wert für mich) einen hohen Wert für mich, während C38 (276 Aufwand, 484 Wert für mich) weniger wichtig ist. Dein Vorschlag mit C10 statt C38 hat einen Aufwand von 2.591, was ausgezeichnet ist. Ich akzeptiere diesen Kompromiss, da er mir C10 gibt und dir A07 und C76 erhält.'}\\ \tt \\ \tt ZUSTIMMUNG: {'B38', 'C00', 'C62', 'C56', 'A07', 'C76', 'C10'} 
	  } 
	   } 
	   } 
	  \\ 
 

    \theutterance \stepcounter{utterance}  

    & & & \multicolumn{2}{p{0.3\linewidth}}{\cellcolor[rgb]{0.95,0.95,0.95}{%
	\makecell[{{p{\linewidth}}}]{% 
	  \tt {\tiny [GM$|$GM]}  
	 game successful 
	  } 
	   } 
	   } 
	 & & \\ 
 

    \theutterance \stepcounter{utterance}  

    & & & \multicolumn{2}{p{0.3\linewidth}}{\cellcolor[rgb]{0.95,0.95,0.95}{%
	\makecell[{{p{\linewidth}}}]{% 
	  \tt {\tiny [GM$|$GM]}  
	 end game 
	  } 
	   } 
	   } 
	 & & \\ 
 

\end{supertabular}
}

\end{document}
